\documentclass[12pt]{article}
\usepackage[margin=1in]{geometry}
\usepackage{amsmath, amssymb, amsthm}
\usepackage{bbm}

\title{\textbf{Advanced Probability Theory}\\
GU 4156: Spring 2025\\
\large Homework \#3}
\author{Professor Ioannis Karatzas\\
Department of Mathematics, Columbia University}
\date{March 12, 2025}

\begin{document}
\maketitle

\noindent This assignment develops an elementary approach to the Second Law of Thermodynamics, in the simple context of a Markov Chain with countable state space. It uses martingale theory, and proceeds via reversal of time.

\bigskip

\noindent\textbf{Exercise \#1:} On a probability space $(\Omega, \mathcal{F}, P)$ consider a Markov Chain $Z = \{Z_n\}_{n \in \mathbb{N}_0}$ with countable state-space $S$, and time-homogeneous transition probabilities given by a matrix $\Pi = \{\pi_{xy}\}_{(x,y) \in S^2}$ with entries
\[
    \pi_{xy} := P(Z_{n+1} = y \mid Z_n = x), \quad \forall\, n \in \mathbb{N}_0.
\]
Denote by $\mathbb{F} = \{\mathcal{F}_n\}_{n \in \mathbb{N}_0}$ the filtration generated by this Chain, to wit, $\mathcal{F}_n := \sigma(Z_0, Z_1, \cdots, Z_n)$.

Suppose also that the Markov Chain $Z = \{Z_n\}_{n \in \mathbb{N}_0}$ is irreducible, positive recurrent, and aperiodic, with initial $P$-distribution $P^{(0)} = \{p(0,x)\}_{x \in S}$ a probability vector with positive components: $p(0,x) = P(Z_0 = x) > 0$ for all $x \in S$.

We saw last semester that such a Chain has a unique invariant distribution: i.e., a vector $Q = \{q(y)\}_{y \in S}$ of positive numbers adding up to 1 and satisfying
\begin{equation}
    q(y) = \sum_{z \in S} q(z)\,\pi_{zy}, \quad \forall\, y \in S. \tag{0.1}
\end{equation}
The $k$-step transition probabilities
\begin{equation}
    \pi^{(0)}_{xy} := \mathbf{1}_{x=y}, \quad \pi^{(k)}_{xy} := P\!\left(Z_k = y \mid Z_0 = x\right), \quad k \in \mathbb{N} \tag{0.2}
\end{equation}
converge as $k$ tends to infinity to $q(y)$, for every pair of states $(x,y) \in S^2$; and the finite-dimensional $P$-distributions of the Chain are given by
\begin{equation}
    P\!\left(Z_0 = x,\, Z_1 = y_1,\, \cdots,\, Z_n = y_n\right) = p(0,x)\,\pi_{xy_1}\,\pi_{y_1 y_2} \cdots \pi_{y_n y_{n+1}}. \tag{0.3}
\end{equation}
Denote also by $P^{(n)} = \{p(n,y)\}_{y \in S}$ the $P$-distribution of the random variable $Z_n$: a column vector with components
\[
    p(n,y) := P(Z_n = y).
\]

Introduce now on the measurable space $(\Omega, \mathcal{F})$ another probability measure $Q$; under this measure, the initial distribution of the Chain is the vector $Q = \{q(y)\}_{y \in S}$ as in (0.1), and thus the finite-dimensional $Q$-distributions of the Chain are given by
\begin{equation}
    Q\!\left(Z_0 = x,\, Z_1 = y_1,\, \cdots,\, Z_n = y_n\right) = q(x)\,\pi_{xy_1}\,\pi_{y_1 y_2} \cdots \pi_{y_n y_{n+1}}. \tag{0.4}
\end{equation}
Then the $Q$-distribution $Q^{(n)} = \{q(n,y)\}_{y \in S}$ of $Z_n$ is $Q^{(n)} \equiv Q$ for all $n \in \mathbb{N}_0$: started at the invariant distribution, the Chain stays with this distribution forever.

\begin{enumerate}
    \item[(i)] Argue that the measure $P$ is absolutely continuous with respect to the measure $Q$ on $\sigma(Z_n)$ for every $n \in \mathbb{N}_0$, with Radon-Nikod\'{y}m derivative (``likelihood ratio'')
    \begin{equation}
        \frac{dP}{dQ}\bigg|_{\sigma(Z_n)} = \ell\!\left(n, Z_n\right), \quad \text{with} \quad \ell(n,y) := \frac{p(n,y)}{q(y)}, \quad (n,y) \in \mathbb{N}_0 \times S. \tag{0.5}
    \end{equation}

    \item[(ii)] Introduce the so-called \emph{relative entropy}
    \[
        H_{P^{(n)}}\!\left(\,\cdot\, \middle\|\, Q\right) := E_P\!\left[\log \ell(n, Z_n)\right]
    \]
    of the distribution $P^{(n)}$ with respect to the distribution $Q$ and argue that, with the function
    \[
        f(u) := u\log(u), \quad u > 0,
    \]
    this quantity can be cast as
    \begin{equation}
        H_{P^{(n)}}\!\left(\,\cdot\, \middle\|\, Q\right) = E_Q\!\left[\log\ell(n,Z_n) \cdot \frac{dP}{dQ}\bigg|_{\sigma(Z_n)}\right] = E_Q\!\left[f\!\left(\ell(n,Z_n)\right)\right]. \tag{0.6}
    \end{equation}

    \item[(iii)] Fix now an integer $T \in \mathbb{N}$ and consider the \emph{time-reversal}
    \begin{equation}
        \hat{Z}_t := Z_{T-t}, \quad t = 0, 1, \cdots, T \tag{0.7}
    \end{equation}
    of the Chain over the finite stretch $0, 1, \cdots, T$; as well as the filtration $\hat{\mathbb{G}} = \{\hat{\mathcal{G}}_t\}_{t=0,1,\cdots,T}$ generated by this time-reversed sequence:
    \[
        \hat{\mathcal{G}}_t = \sigma\!\left(\hat{Z}_0, \hat{Z}_1, \cdots, \hat{Z}_t\right) = \sigma\!\left(Z_T, Z_{T-1}, \cdots, Z_{T-t}\right).
    \]
    Argue that, under each one of the probability measures $P$ and $Q$, the collection of random variables in (0.7) has the Markov property. What are the corresponding one-day transition probabilities? Are they time-homogeneous?

    \bigskip
    \noindent\textit{Assume from now on$^1$ also that the state-space $S$ is finite.}

    \item[(iv)] With the notation of (0.5), (0.7), show that
    \[
        \left\{\ell\!\left(T-t,\, \hat{Z}_t\right)\right\}_{t = 0, 1, \cdots, T}
    \]
    is a $(\hat{\mathbb{G}}, Q)$-martingale.

    \item[(v)] Show that the sequence $\mathbb{N}_0 \ni n \mapsto H_{P^{(n)}}\!\left(\,\cdot\,\middle\|\,Q\right) \in [0,\infty)$ of relative entropies across time is nonnegative, decreasing, and dissipates down to zero, in the sense
    \[
        \lim_{n \to \infty} \downarrow H_{P^{(n)}}\!\left(\,\cdot\,\middle\|\,Q\right) = 0.
    \]

    \item[(vi)] Argue that, for the Total Variation distance
    \[
        \left\|\mu - \nu\right\|_{\mathrm{TV}} := \sup_{A \in \mathcal{F}}\left|\mu(A) - \nu(A)\right|
    \]
    between probability measures, we have also
    \[
        \lim_{n \to \infty} \left\|P^{(n)} - Q\right\|_{\mathrm{TV}} = 0.
    \]
\end{enumerate}

\bigskip

\noindent\textbf{Remark:} The results in (v), (vi) express of course the Second Law of Thermodynamics in the present context.

\bigskip

\noindent\textbf{Remark:} Nothing about the function $f(u) = u\log(u)$, $u > 0$, in (0.6), is exactly ``sacrosanct'' in (iv), (v). We can replace it by any convex function $F : (0,\infty) \to [0,\infty)$ with $F(1) = 0$,$^2$ define the so-called ``$F$-relative entropy''
\[
    H^F\!\left(P^{(n)}\,\middle\|\,Q\right) := E_P\!\left[F\!\left(\frac{\ell(n,Z_n)}{\ell(n,Z_n)}\right)\right] = E_Q\!\left[F\!\left(\ell(n,Z_n)\right)\right]
\]
in the manner of (ii) above; then everything goes through just as well.

On the other hand, (vi) does need the special form of the original relative entropy with the logarithmic function.

\vspace{1cm}
\noindent\rule{0.5\textwidth}{0.4pt}

\noindent $^1$ For simplicity only. The results (iv)--(vi) below continue to hold for countable state-spaces, under the condition $H_{P^{(0)}}(\,\cdot\,\|\,Q) < \infty$. Give this a try as well.

\noindent $^2$ For instance, $F(u) = u^m - 1$, $u \in (0,\infty)$, for any $m > 1$.

\end{document}
